% Section skeleton for the paper (Requirement 4.5).
% Open the SIAM template in Overleaf:
%   https://www.overleaf.com/latex/templates/template-for-siam-journals/cbbdwqgfhrkd
% click "Open as Template", then paste these sections into the document body,
% replacing the template's example sections. Keep them in this order.
% Target: about 10 pages, and no fewer than 8 excluding references (Requirement 4.3).
% Export the compiled PDF as paper.pdf at the root of this repository.

% The abstract goes in the template's \begin{abstract} ... \end{abstract} block.
% TODO: abstract

\section{Introduction}
% TODO

\section{Research question and relevant literature}
% TODO

\section{Methodology}
% TODO

\section{Data set: source and cleaning}
% TODO: where the data came from, its licence, and every cleaning step (see data/README.md).

\section{Implementation}
% TODO: how the code in code/ is organised and run.

\section{Results}
% TODO: every figure and table here must be produced by code/04_figures.ipynb (Requirement 4.11).
% Example:
% \begin{figure}[htbp]
%   \centering
%   \includegraphics[width=0.8\linewidth]{fig1_name.pdf}
%   \caption{What the reader should look at.}
%   \label{fig:fig1}
% \end{figure}

\section{Conclusion}
% TODO

% References: \bibliographystyle{siamplain} \bibliography{references}

\appendix
\section{Declaration of tools used}
% Requirements 5 and 6: list every generative AI tool / code assistant and what it was used for,
% and every piece of code you did not write, with its source and licence.
\begin{itemize}
  \item Claude Code: generated the initial repository scaffold (folder layout, \texttt{code/helpers.py},
        notebook templates, README template and this section outline).
  \item TODO: other tools, and what each was used for.
\end{itemize}

\subsection*{Reused code}
% TODO: source, link, where it is used, licence. Course material: identify where used.
None so far.
